\chapter*{读完后的检查表}

读完本书后，至少应该能回答下面这些问题：

\begin{longtable}{p{0.25\linewidth}p{0.65\linewidth}}
\toprule
主题 & 能力要求 \\
\midrule
比特 & 能说明为什么计算机用两个稳定区间表示 0 和 1，而不是追求连续电压的精确值。 \\
逻辑门 & 能从开关路径解释 NOT、AND、OR、NAND 的行为，并理解为什么 NAND 可以构造其他门。 \\
组合逻辑 & 能把真值表转成门电路，解释半加器、全加器、多路选择器和译码器分别解决什么问题。 \\
时序逻辑 & 能说明锁存器、触发器、寄存器和时钟边界的关系，知道为什么组合逻辑不能保存历史。 \\
状态机 & 能把状态、输入、转移条件和输出分开，并说明状态机怎样控制一组硬件动作。 \\
ALU & 能把加法、减法、按位运算、比较和标志位放在同一个计算部件里解释。 \\
最小 CPU & 能画出 PC、指令存储器、寄存器堆、ALU、数据存储器和控制器之间的数据流。 \\
整机硬件 & 能说明主存、总线、MMIO、DMA、中断控制器和复位启动怎样把最小 CPU 扩展成整机。 \\
平台事务 & 能把一次内存、MMIO、DMA 或中断访问写成发起者、目标、地址、握手、完成、错误和副作用的 trace。 \\
\bottomrule
\end{longtable}

\begin{classicnote}
最终检查不以“能不能复述章节标题”为准，而以能不能给出机器证据为准。

\begin{itemize}
\item 给一个信号节点，能说出它何时被拉高、何时被拉低、何时悬空、何时冲突。
\item 给一个组合电路，能列真值表、最长路径、可能毛刺和后级采样边界。
\item 给一个状态电路，能指出旧状态、下一状态、时钟沿、建立/保持约束和复位行为。
\item 给一个算术结果，能分别按无符号和补码解释，并说明 Carry 与 Overflow 的差别。
\item 给一条最小 CPU 指令，能写出控制信号 trace、数据通路 trace 和 PC 下一值。
\item 给一次内存或设备访问，能写出地址译码、总线/互连事务、响应返回和完成证据。
\item 给一次 DMA 或中断，能写出缓冲区所有权、可见性、状态清除和完成通知顺序。
\end{itemize}

如果这些证据能写出来，说明读者已经把“电平 -> 门 -> 组合逻辑 -> 状态 -> 数据通路 -> 控制器 -> CPU -> 整机”的链路连成了一台机器。
\end{classicnote}
